\section{信源编码：用最短的平均码长把消息写下来}

\begin{idea}{这一讲在解决什么问题（先看地图）}
前几讲是\textbf{算账}：用熵 $\Hx{X}$ 量出一个信源到底有多少不确定性。这一讲开始干活：
硬盘和网线只认 $0$ 和 $1$，所以得把信源符号\textbf{换成一串 0/1}，这叫\textbf{编码}。
换法有无穷多种，凭什么说一种比另一种好？两条标准：
\textbf{(1) 能还原}——收到 0/1 串后必须一模一样地恢复原文（本章叫"无错编码"）；
\textbf{(2) 在能还原的前提下尽量省}——平均每个符号花的比特数尽量少。
于是全讲主线只有一句话：\textbf{在"能唯一还原"的码里把平均码长压到最小，而这条平均码长的下限就是熵。}
码的分类（什么叫能还原）$\to$ Kraft 不等式（什么码长可行）$\to$ 两个界限（熵是下限）
$\to$ Huffman（怎么把最小值做出来）$\to$ 熵率（有前后文时下限变成什么）。
\end{idea}

%==============================================================
\subsection{编码到底在干什么：把符号换成 0/1 串}

存一篇文章进硬盘，硬盘里没有"字"，只有能存高/低电平的格子，即只认 $0$ 和 $1$。
所以必须事先约定一张\textbf{对照表}：这个字对应这串 0/1，那个字对应那串 0/1；
存的时候查表写，读的时候反查表还原。最熟的例子是 ASCII 码：字母 \texttt{A} 记作 $1000001$（7 位）。

生活比方：\textbf{寄快递}。信源符号形状各异，不能直接扔上传送带，得先装进标准纸箱（码字）。
箱子有大有小（码字有长有短），\textbf{装得越紧凑，同样的车皮运得越多}——这就是"平均码长越小越好"。

\begin{defbox}{五个名词（定义 4.3.2--4.3.4）}
设信源符号集 $S=\{s_1,\dots,s_q\}$，$s_i$ 出现的概率为 $p_i=\p{s_i}$，$\sum_i p_i=1$；
码符号表（字母表）为 $A=\{a_1,\dots,a_r\}$。
\begin{itemize}[leftmargin=1.3em,itemsep=1pt,topsep=2pt]
  \item $r$ 叫码的\textbf{基数}：只用 $0,1$ 时 $r=2$（\textbf{二元码}），用 $0,1,2$ 时 $r=3$（\textbf{三元码}）。
  \item \textbf{码字} $w_i$：分给 $s_i$ 的那串码符号，如 $w_2=0110$。
  \item \textbf{码长} $l_i=\mathrm{len}(w_i)$：码字里码符号的个数，上例 $l_2=4$。
  \item \textbf{码} $C=\{w_1,\dots,w_q\}$：整张对照表。
  \item \textbf{编码规则} $(C,f)$，$f:S\to C$：把每个符号送到它的码字，
        且必须是\textbf{单射}（不同符号不能共用一个码字），否则一开始就没法还原。
\end{itemize}
\end{defbox}

\begin{defbox}{平均码长（本讲第一主角，必背）}
\begin{align*}
\bar L=\mathrm{Avelen}(C,f)=\sum_{i=1}^{q}p_i\,l_i=p_1l_1+p_2l_2+\cdots+p_ql_q .
\end{align*}
人话：\textbf{把各码字的长度按它出现的频繁程度加权平均}。出现得多的符号，
它的码长在平均值里说话更有分量。所以想把 $\bar L$ 压低，最有效的办法就是
\textbf{给常出现的符号配短码字}。注意 $\bar L$ 只与"概率 + 码长"有关，
跟码字具体长什么样（是 $0110$ 还是 $1011$）无关。
\end{defbox}

\begin{calc}{手算一个三符号的平均码长（不跳步）}
$p_1=0.5,\ p_2=0.3,\ p_3=0.2$（和 $=1$），二元码如下表。

\begin{center}
\begin{tabular}{@{}cccc@{}}
\toprule
符号 & 概率 $p_i$ & 码字 $w_i$ & 码长 $l_i$\\
\midrule
$s_1$ & $0.5$ & $0$  & $1$\\
$s_2$ & $0.3$ & $10$ & $2$\\
$s_3$ & $0.2$ & $11$ & $2$\\
\bottomrule
\end{tabular}
\end{center}

\textbf{第一步抄公式}：$\bar L=\sum_i p_il_i$。\textbf{第二步逐项相乘}：
$\bar L=0.5\times1+0.3\times2+0.2\times2=0.5+0.6+0.4=1.5$ bit/符号。
\textbf{第三步跟熵比}（用 $\Hx{X}=\sum_i p_i\lb\frac{1}{p_i}$，先算每个 $\lb\frac1{p_i}$）：
\begin{align*}
\lb\tfrac{1}{0.5}=\lb 2=1,\quad
\lb\tfrac{1}{0.3}=\lb 3.333\approx1.7370,\quad
\lb\tfrac{1}{0.2}=\lb 5\approx2.3219,
\end{align*}
$\Hx{X}=0.5\times1+0.3\times1.7370+0.2\times2.3219=0.5+0.5211+0.4644=1.4855$ bit/符号。
\textbf{第四步效率}：$\eta=\dfrac{\Hx{X}}{\bar L}=\dfrac{1.4855}{1.5}\approx99.03\%$，
即平均每符号只比理论下限多花 $1.5-1.4855=0.0145$ 比特。
\end{calc}

%==============================================================
\subsection{码的三层分类：非奇异 $\subset$ 唯一可译 $\subset$ 前缀码}

\subsubsection{先讲你天天遇到的问题：电话号码}
假设公司分机号是 $1$、$12$、$123$。你拨 "$12$"，交换机听见 "1" 就以为你要转 $1$ 号分机，
"2" 被当成下一个动作——\textbf{永远拨不到 $12$ 号}。这就是"码字不能互为前缀"的意思：
\textbf{短号码不能是长号码的开头}，否则一听到短号码就以为该停了，后面几位没法解释。

\begin{defbox}{三层定义（4.3.5--4.3.6，必背）}
\textbf{(1) 非奇异码}：$f$ 是单射，不同符号给不同码字。
人话：\textbf{一张表里没有两个相同的码字}。\\
\textbf{(2) 唯一可译码}：任一串码字 $c_{i_1}\cdots c_{i_k}$ 的分割方式唯一。
人话：\textbf{收到长 0/1 串后能一刀一刀切回原文，且切法只有一种}。\\
\textbf{(3) 前缀码（异字头码、即时码）}：$C$ 中没有任何码字是另一个码字的前缀。
人话：\textbf{一个码字一收完就知道它结束了}，可以立刻译出，不用等后面（"即时"、"无译码延时"）。
\end{defbox}

三者\textbf{层层包含}（选择题最爱考）：
\begin{align*}
\underbrace{\text{前缀码（即时码）}}_{\text{最强，无延时}}\;\subsetneq\;
\underbrace{\text{唯一可译码}}_{\text{能还原，可能要等}}\;\subsetneq\;
\underbrace{\text{非奇异码}}_{\text{最弱，只是不重样}} .
\end{align*}
一句话：\textbf{前缀码一定唯一可译且无延时；唯一可译码不一定是前缀码。}

\subsubsection{自己动手判四张码表（A/B/C/D 四型）}
同一信源、概率是"折半"型 $p=(0.5,\,0.25,\,0.125,\,0.125)$（和 $=1$），四张二元码表：

\begin{center}
\begin{tabular}{@{}ccccc|c@{}}
\toprule
符号 & 概率 $p_i$ & 码 A & 码 B & 码 C & 码 D\\
\midrule
$s_1$ & $0.5$   & $0$  & $0$  & $0$    & $0$\\
$s_2$ & $0.25$  & $1$  & $1$  & $01$   & $10$\\
$s_3$ & $0.125$ & $1$  & $00$ & $011$  & $110$\\
$s_4$ & $0.125$ & $1$  & $11$ & $0111$ & $111$\\
\midrule
平均码长 $\bar L$ & & $1.0$ & $1.25$ & $1.875$ & $1.75$\\
\bottomrule
\end{tabular}
\end{center}

\textbf{码 A}（$0,1,1,1$）：$s_2,s_3,s_4$ 共用码字 "$1$"，收到 "$1$" 分不清是哪个，
连\textbf{非奇异码}都不是。$\bar L_A=0.5+0.25+0.125+0.125=1$。\\
\textbf{码 B}（$0,1,00,11$）：四个码字互不相同，\textbf{非奇异} $\checkmark$。
但 "$11$" 既可切成 $s_4$、也可切成 $s_2s_2$，两种切法都合法，
故\textbf{不是唯一可译码}。$\bar L_B=0.5\times1+0.25\times1+0.125\times2+0.125\times2=1.25$。\\
\textbf{码 C}（$0,01,011,0111$）："$0$" 是 "$01$" 的前缀，故\textbf{不是前缀码}；
但它\textbf{是唯一可译的}——四个码字都只有一个开头的 $0$，后面 $1$ 的连续个数决定是哪个符号
（$0$ 个 $1$ 是 $s_1$，$1$ 个是 $s_2$，$2$ 个是 $s_3$，$3$ 个是 $s_4$）。
麻烦在于收到 "$011$" 时不能确定它是 $s_3$ 到头了、还是 $s_4=0111$ 的前三位，
\textbf{必须等下一个符号才敢拍板}——这就是\textbf{有译码延时}。
$\bar L_C=0.5\times1+0.25\times2+0.125\times3+0.125\times4=1.875$。\\
\textbf{码 D}（$0,10,110,111$）：$0$ 开头一个、$1$ 后接 $0/10/11$ 各不相同，无前缀关系，
是\textbf{前缀码}，收完即译。$\bar L_D=0.5\times1+0.25\times2+0.125\times3+0.125\times3=1.75$。

\begin{idea}{平均码长 $=1$ 的那个码（码 A）是不是最好的？不是，而且不能用}
表里 $\bar L$ 最小的恰是码 A（$=1$），只看数字 A 好像最好。但这个"最好"是\textbf{假的最好}：
\begin{itemize}[leftmargin=1.3em,itemsep=1pt,topsep=2pt]
  \item \textbf{它根本不能用}：三个符号共用 "$1$"，收到无法判断原文，连非奇异码都不满足。
        还原不回来的码，$\bar L$ 再小也没人要。
  \item \textbf{平均码长只在"唯一可译码"这一族里比才有意义}：把不能还原的码也算进来，
        $\bar L$ 可以任意小（所有符号都编成空串就 $\bar L=0$），但那等于什么都没写，信息全丢了。
\end{itemize}
\textbf{正确的问法是}：在唯一可译码（工程上用前缀码）这一族里，$\bar L$ 最小的是谁？
本表唯一可译的只有 C 和 D，D 更小（$1.75<1.875$）且无延时，
所以\textbf{四张表里最好的是 D}，而且它恰好达到熵（见下），一分钱没浪费。
\end{idea}

\textbf{顺手验证 D 为什么能到 $100\%$}：这个信源概率都是 $2$ 的负整数次幂，
\begin{align*}
\Hx{X}=0.5\lb 2+0.25\lb 4+0.125\lb 8+0.125\lb 8
=0.5+0.5+0.375+0.375=1.75=\bar L_D .
\end{align*}
效率 $\eta=100\%$。规律：当 $p_i=2^{-k_i}$ 时最优码长就是 $l_i=\lb\frac{1}{p_i}$（整数），
此时 $\bar L=\sum p_i\lb\frac1{p_i}=\Hx{X}$，一分冗余都没有。

%==============================================================
\subsection{Kraft 不等式：码字空间到底够不够用}

\subsubsection{人话版：一分四、四分十六}
拿一张纸代表"所有可能的 0/1 串"，只能靠\textbf{二分}来划分：中间剪一刀，左边串以 $0$ 开头、
右边以 $1$ 开头。于是\textbf{长度 1 的码字（如 "$0$"）占掉半张纸}（份额 $1/2$），
\textbf{长度 2 的码字（如 "$01$"）只占 $1/4$，长度 $l$ 的占 $\frac{1}{2^l}$}——
因为长度 $l$ 的二进制串共有 $2^l$ 个，挑走一个就是切走 $\frac{1}{2^l}$；$r$ 元码同理占 $\frac{1}{r^l}$。

\textbf{码字之间不能互相盖住}（某个码字一旦成了别人的前缀，就把别人那块纸整个吞了），
而纸总共只有 $1$ 张。所以\textbf{所有码字占的份额加起来不能超过 $1$}——
这就是 Kraft 不等式，一点也不神秘。

\begin{defbox}{Kraft 不等式（定理 4.3.1，必背）}
设码字长度为 $l_1,\dots,l_q$，则
\begin{align*}
\sum_{i=1}^{q}r^{-l_i}=r^{-l_1}+r^{-l_2}+\cdots+r^{-l_q}\ \le\ 1 .
\end{align*}
\textbf{必要性}：若 $C$ 是 $r$ 元即时码（前缀码），上式必成立——\textbf{真编得出来就一定没超支}。\\
\textbf{充分性}：若一组正整数 $l_1,\dots,l_q$ 满足上式，则\textbf{一定存在}具有这些码长的
$r$ 元前缀码——\textbf{没超支就一定编得出来}。\\
并且 McMillan 定理保证：对\textbf{唯一可译码}该不等式同样成立，所以这套判据对唯一可译码也管用。
\end{defbox}

\textbf{一句话翻译}：$\sum_i r^{-l_i}\le1$ 说的是"你打算用的这些码长，占的位置\textbf{不超过}整个码字空间"。
长码字占得少（$\frac{1}{r^l}$ 小），短码字占得多（$\frac{1}{r^l}$ 大），
所以"用太多短码字"必然超支——这正是"不可能让所有符号都用很短的码字"的原因。

\begin{calc}{判断二元/三元前缀码的存在性（考试常见小问）}
\textbf{(1) 码长 $\{1,2,3,3\}$ 能否构成二元前缀码？}（$r=2$，逐项算 $2^{-l}$）
\begin{align*}
\sum_i2^{-l_i}=\frac12+\frac14+\frac18+\frac18
=\frac{4}{8}+\frac{2}{8}+\frac{1}{8}+\frac{1}{8}=\frac{8}{8}=1\le1 .
\end{align*}
\textbf{能。}和恰为 $1$，说明纸被切得一点不剩，叫\textbf{满树码（完全码）}。
具体例子就是上面的码 D：$0,10,110,111$。

\textbf{(2) 码长 $\{1,2,2,3\}$ 能否？}
\begin{align*}
\sum_i2^{-l_i}=\frac12+\frac14+\frac14+\frac18
=\frac{4}{8}+\frac{2}{8}+\frac{2}{8}+\frac{1}{8}=\frac{9}{8}=1.125>1 .
\end{align*}
\textbf{不能。}超支 $1/8$ 张纸。直观：两个长度 $2$ 的码字吃掉 $1/2$，再加一个长度 $1$ 的又是 $1/2$，
这三个就占满 $1$ 了，第四个（哪怕长度 $3$）无处安放。

\textbf{(3) 三元码、长度 $\{1,2,2,2,2,2,3,3\}$}（$r=3$）：
\begin{align*}
\sum_i3^{-l_i}=\frac13+5\times\frac19+2\times\frac{1}{27}
=\frac{9}{27}+\frac{15}{27}+\frac{2}{27}=\frac{26}{27}\approx0.963\le1 .
\end{align*}
成立，故存在这样的三元前缀码（下面三元 Huffman 的例子正好给出它）。注意 $0.963<1$，
还剩一点空间没用——只有概率都是 $3$ 的负整数次幂时才可能刚好填满。
\end{calc}

\textbf{怎么用它判断"某码长分配可行不可行"}（三步走）：
(1) 数清是几元码，定下 $r$；
(2) 把每个 $\frac{1}{r^{l_i}}$ 算出来相加（\textbf{通分}成同一分母最不易错，上例都化成 $/8$、$/27$）；
(3) 和 $\le1$ $\Rightarrow$ \textbf{可行}（存在这样的前缀码）；和 $>1$ $\Rightarrow$ \textbf{不可行}；
若和 $=1$，补一句"满树码"，老师会加分。

\begin{qa}{Kraft 不等式成立，是不是就说明"这张码表"是前缀码？}
\textbf{不是，这是最常见的误解。} Kraft 定理管的是\textbf{"码长这组数字"}：它保证存在某个前缀码
其码长恰是这些数，但\textbf{不保证你手上这张具体码表}是前缀码。
反例（教材例 4.3.3）：$C=\{0,\,11,\,100,\,1100\}$，长度 $1,2,3,4$，
$2^{-1}+2^{-2}+2^{-3}+2^{-4}=\frac12+\frac14+\frac18+\frac1{16}=\frac{15}{16}\le1$
完全成立，但 "$11$" 是 "$1100$" 的前缀，所以这个码\textbf{不是即时码}。
\textbf{分工要记清}：判断一组\textbf{长度数字}用 Kraft；判断一张\textbf{具体码表}就逐个检查"谁是谁的前缀"。
\end{qa}

%==============================================================
\subsection{变长编码的两个界限：熵是理论下限}

\begin{defbox}{无失真变长编码定理（Shannon 第一定理，必背）}
对离散无记忆信源 $X$ 作 $r$ 元变长编码、平均码长 $\bar L$，则\textbf{任何}唯一可译码都满足
$\bar L\ge\dfrac{\Hx{X}}{\lb r}$，而且\textbf{一定存在}唯一可译码使其满足
$\bar L<\dfrac{\Hx{X}}{\lb r}+1$。二元码（$r=2$，$\lb2=1$）就是最常用的版本：
\begin{align*}
\boxed{\ \Hx{X}\ \le\ \bar L\ <\ \Hx{X}+1\ }
\end{align*}
\end{defbox}

\textbf{逐句翻译}：
\textbf{左边 $\Hx{X}\le\bar L$} 是\textbf{硬下限，谁都突破不了}。直观理由：$N$ 个符号的消息共 $q^N$ 种组合，
要区分它们至少需要约 $\lb q^N$ 比特，而熵正是"平均到底需要多少比特"的精确答案。
\textbf{右边 $\bar L<\Hx{X}+1$} 说\textbf{好码一定存在，且最多只比熵多不到 $1$ 比特}：
多出的那点可理解为"取整的代价"，因为码长必须是整数，而 $-\lb p_i$ 一般是小数，向上取整最多浪费 $1$ 位。
合起来：$\bar L$ 被夹在宽度为 $1$ 的窗口 $[\Hx{X},\Hx{X}+1)$ 里；
想再逼近就得\textbf{分块编码}（把 $N$ 个符号当一个符号编，界限变成
$\Hx{X}\le\frac{\bar L_N}{N}<\Hx{X}+\frac1N$，$N$ 越大越贴近熵）。

\textbf{"为什么能这么近"的构造思路}（一句话版）：直接把码长取成
$l_i=\left\lceil\lb\frac{1}{p_i}\right\rceil$（$\lceil\cdot\rceil$ 是\textbf{向上取整}，
如 $\lceil2.32\rceil=3$）。于是 $\lb\frac1{p_i}\le l_i<\lb\frac1{p_i}+1$，推出
$2^{-l_i}\le p_i$，求和得 $\sum_i2^{-l_i}\le\sum_ip_i=1$，由 Kraft 知这样的前缀码存在；
而 $\bar L=\sum_ip_il_i<\sum_ip_i\left(\lb\frac1{p_i}+1\right)=\Hx{X}+1$。
\textbf{右边的不等号就是这么来的：取整最多浪费 1 位。}

\begin{calc}{$H=2.6087$、$\bar L=2.72$ 的码到底有多好}
拿本节后面 Huffman 那个 $7$ 符号例子的结果核对：
$\Hx{X}=2.6087$ bit/符号，$\bar L=2.72$ bit/符号。
\textbf{先验界限}：$2.6087\le2.72<2.6087+1=3.6087$，两边都满足 $\checkmark$。
\textbf{再算效率}（实际离下限有多近）：
$\eta=\dfrac{\Hx{X}}{\bar L}=\dfrac{2.6087}{2.72}\approx0.9591=95.91\%$。
\textbf{再算冗余度}（浪费了多少）：$1-\eta=1-0.9591=0.0409=4.09\%$，
即每符号多花 $2.72-2.6087=0.1113$ 比特。
\end{calc}

读法：$\bar L$ 只比熵多 $0.111$ 比特，远小于定理允许的 $1$ 比特，说明 Huffman 码已非常接近极限。
考试凡让你"算平均码长并评价"，就要把 $\Hx{X}$、$\bar L$、$\eta$、冗余度\textbf{四个数都写出来}，
再补一句"$\Hx{X}\le\bar L<\Hx{X}+1$ 成立"。

%==============================================================
\subsection{Huffman 编码：把最小平均码长真的做出来}

\subsubsection{一句话原理：每次把最轻的两个箱子捆一起}
还是寄快递：$7$ 个包裹重量（概率）不同，要装成一个大包裹，路径用 $0/1$ 标记。
显然\textbf{越重的包裹该越早被分拣出去（路径越短）}；
而\textbf{最轻的两个最该被"捆在一起"共享同一段路径}——它们最少出现，让它们多走几步代价最小。
做法：\textbf{反复把概率最小的两个合并成新节点（概率相加），直到只剩一个节点（概率 $=1$）；
再从根往回走，左枝记 $0$、右枝记 $1$，路径上的 $0/1$ 序列就是码字。}

\begin{idea}{四条铁律（记住这四条就不会算错）}
(1) 每次取\textbf{当前序列中最小的两个}相加，把和\textbf{插回}序列并\textbf{重新排序}；
(2) 二元码合并次数固定为 $n-1$ 次；
(3) 码符号分配方向可任选（左 $0$ 右 $1$ 或反之），\textbf{所以 Huffman 码不唯一}；
(4) \textbf{必须画码树或写完整合并过程}，只给最终码表老师看不出你是怎么来的。
\end{idea}

\subsubsection{手把手全过程：7 个符号的经典例子}

\begin{calc}{完整 Huffman 编码（考试最可能的计算题，照着抄一遍）}
概率（已降序）$P=(0.20,\,0.19,\,0.18,\,0.17,\,0.15,\,0.10,\,0.01)$，和 $=1$，
对应符号 $s_1,\dots,s_7$。

\textbf{第一步：合并过程}（老师要看的"过程"，一行都不能省）：

\begin{center}
\begin{tabular}{@{}clc@{}}
\toprule
轮次 & 当前概率序列（从大到小） & 本次合并最小的两个\\
\midrule
初始 & $0.20,\ 0.19,\ 0.18,\ 0.17,\ 0.15,\ 0.10,\ 0.01$ & ---\\
第 1 轮 & $0.20,\ 0.19,\ 0.18,\ 0.17,\ 0.15,\ \mathbf{0.11}$ & $0.10+0.01=0.11$\\
第 2 轮 & $0.26,\ 0.20,\ 0.19,\ 0.18,\ \mathbf{0.17}$ & $0.15+0.11=0.26$\\
第 3 轮 & $0.35,\ 0.26,\ 0.20,\ \mathbf{0.19}$ & $0.17+0.18=0.35$\\
第 4 轮 & $0.39,\ 0.35,\ \mathbf{0.26}$ & $0.19+0.20=0.39$\\
第 5 轮 & $0.61,\ \mathbf{0.39}$ & $0.26+0.35=0.61$\\
第 6 轮 & $\mathbf{1.00}$ & $0.39+0.61=1.00$\\
\bottomrule
\end{tabular}
\end{center}

注意每轮都只取"当前最小的两个"，合并后\textbf{重新按大小插回}；
合并次数 $6=n-1$ 次，最后只剩根节点 $1.00$ $\checkmark$。

\textbf{第二步：倒着展开，分配码字}（从根往下，左枝 $0$、右枝 $1$）：
\begin{align*}
1.00&=\underbrace{0.61}_{0}+\underbrace{0.39}_{1}, &
0.61&=\underbrace{0.35}_{00}+\underbrace{0.26}_{01},\\
0.39&=\underbrace{0.19}_{10}+\underbrace{0.20}_{11}, &
0.35&=\underbrace{0.17}_{000}+\underbrace{0.18}_{001},\\
0.26&=\underbrace{0.15}_{010}+\underbrace{0.11}_{011}, &
0.11&=\underbrace{0.10}_{0110}+\underbrace{0.01}_{0111}.
\end{align*}

\textbf{第三步：读出码表}（顺着下标就能读出码字）：

\begin{center}
\begin{tabular}{@{}ccccc@{}}
\toprule
符号 & 概率 $p_i$ & 码字 $w_i$ & 码长 $l_i$ & $p_il_i$\\
\midrule
$s_1$ & $0.20$ & $11$   & $2$ & $0.40$\\
$s_2$ & $0.19$ & $10$   & $2$ & $0.38$\\
$s_3$ & $0.18$ & $001$  & $3$ & $0.54$\\
$s_4$ & $0.17$ & $000$  & $3$ & $0.51$\\
$s_5$ & $0.15$ & $010$  & $3$ & $0.45$\\
$s_6$ & $0.10$ & $0110$ & $4$ & $0.40$\\
$s_7$ & $0.01$ & $0111$ & $4$ & $0.04$\\
\midrule
合计 & $1.00$ & & $\bar L=$ & $2.72$\\
\bottomrule
\end{tabular}
\end{center}

码长向量 $(2,2,3,3,3,4,4)$——\textbf{概率越大码越短}，符合"大配短、小配长"。
\textbf{顺手用 Kraft 验算}（能立刻发现算错）：
\begin{align*}
\sum_i2^{-l_i}=2\times2^{-2}+3\times2^{-3}+2\times2^{-4}
=\frac{2}{4}+\frac{3}{8}+\frac{2}{16}
=\frac{8}{16}+\frac{6}{16}+\frac{2}{16}=1\ \checkmark
\end{align*}
和恰为 $1$，是一棵\textbf{满树}、没有浪费位置（这是 Huffman 码的特征：
只要所有概率都参与合并，树一定是满的）。
\end{calc}

\begin{calc}{四个数一个不能少（考试得分点）}
\textbf{(1) 平均码长}（上表最后一列相加）：
\begin{align*}
\bar L&=0.20\times2+0.19\times2+0.18\times3+0.17\times3+0.15\times3+0.10\times4+0.01\times4\\
&=0.40+0.38+0.54+0.51+0.45+0.40+0.04=2.72\ \text{bit/符号}.
\end{align*}
\textbf{(2) 熵}（用 $\Hx{X}=\sum_ip_i\lb\frac1{p_i}$，先把每个 $\lb\frac1{p_i}$ 单独算出来）：
\begin{align*}
\lb\tfrac{1}{0.20}&=\lb5\approx2.3219, &
\lb\tfrac{1}{0.19}&\approx2.3959, &
\lb\tfrac{1}{0.18}&\approx2.4739,\\
\lb\tfrac{1}{0.17}&\approx2.5564, &
\lb\tfrac{1}{0.15}&\approx2.7370, &
\lb\tfrac{1}{0.10}&=\lb10\approx3.3219,\\
\lb\tfrac{1}{0.01}&=\lb100\approx6.6439 . &&&
\end{align*}
（校验：$2^{2.3219}\approx5$，$2^{6.6439}\approx100$，对得上。）逐项相乘再求和：
\begin{align*}
\Hx{X}&=0.20\times2.3219+0.19\times2.3959+0.18\times2.4739+0.17\times2.5564\\
&\quad+0.15\times2.7370+0.10\times3.3219+0.01\times6.6439\\
&=0.4644+0.4552+0.4453+0.4346+0.4105+0.3322+0.0664\\
&=2.6087\ \text{bit/符号}.
\end{align*}
（中间相加：$0.4644+0.4552=0.9196$，$+0.4453=1.3649$，$+0.4346=1.7995$，
$+0.4105=2.2100$，$+0.3322=2.5422$，$+0.0664=2.6086$，末位舍入差 $0.0001$。）

\textbf{(3) 效率}：$\eta=\dfrac{\Hx{X}}{\bar L}=\dfrac{2.6087}{2.72}\approx0.9591=95.91\%$。\\
\textbf{(4) 冗余度}：$1-\eta=1-0.9591=0.0409=4.09\%$。\\
\textbf{(5) 对照定理}：$2.6087=\Hx{X}\le\bar L=2.72<3.6087=\Hx{X}+1$ $\checkmark$，
且只比下限多 $0.1113$ 比特，比定理允许的 $1$ 比特好得多。
\end{calc}

\begin{idea}{为什么 Huffman 码最优（一句话版）}
教材定理 4.3.3：Huffman 码在所有 $r$ 元即时码中\textbf{平均码长最小}；
再由"所有唯一可译码的最小平均码长与所有即时码的最小平均码长\textbf{相同}"（推论 4.3.2），
知它在所有唯一可译码里也最优。考试一般不要求证明，但要求\textbf{会用}。
\end{idea}

\subsubsection{三元 Huffman 编码：首步要合并几个？}
二元码每次合 $2$ 个，一路下去总会剩 $1$ 个；但 $r$ 元码每次合 $r$ 个，
\textbf{可能最后剩下的不是 $1$ 个}，树就长不完整。解决办法：\textbf{第一次先少合几个}。

\begin{defbox}{首步合并个数（必背公式）}
对 $n$ 个符号作 $r$ 元 Huffman 编码，\textbf{第一轮合并}
\begin{align*}
s=1+\big[(n-1)\bmod(r-1)\big]
\end{align*}
个（若算出 $s=1$，说明不用调整，每轮都合 $r$ 个即可），之后\textbf{每轮都合 $r$ 个}。
人话：\textbf{要让每轮把"待合并个数"减少 $r-1$、最后正好减到 $1$，
就得先用一次小合并把个数调成 $1\bmod(r-1)$ 的样子。}
\end{defbox}

\begin{calc}{三元 Huffman 的小例子（$n=8$，$r=3$）}
概率（降序）$0.22,0.20,0.18,0.15,0.10,0.08,0.05,0.02$（和 $=1$）。
\textbf{先算首步个数}：$n=8$、$r=3$、$r-1=2$，
$s=1+[(8-1)\bmod2]=1+[7\bmod2]=1+1=2$，
所以\textbf{第一轮只合 $2$ 个}，之后每轮合 $3$ 个。走一遍：

\begin{center}
\small
\begin{tabular}{@{}clc@{}}
\toprule
轮次 & 当前概率序列（从大到小） & 本次合并\\
\midrule
初始 & $0.22,\ 0.20,\ 0.18,\ 0.15,\ 0.10,\ 0.08,\ 0.05,\ 0.02$ & ---\\
第 1 轮（合 $2$ 个） & $0.22,\ 0.20,\ 0.18,\ 0.15,\ 0.10,\ 0.08,\ \mathbf{0.07}$ & $0.05+0.02=0.07$\\
第 2 轮（合 $3$ 个） & $0.25,\ 0.22,\ 0.20,\ 0.18,\ \mathbf{0.15}$ & $0.10+0.08+0.07=0.25$\\
第 3 轮（合 $3$ 个） & $0.53,\ 0.25,\ \mathbf{0.22}$ & $0.15+0.18+0.20=0.53$\\
第 4 轮（合 $3$ 个） & $\mathbf{1.00}$ & $0.22+0.25+0.53=1.00$\\
\bottomrule
\end{tabular}
\end{center}

\textbf{分配码字}（三元码用 $0,1,2$）：
\begin{align*}
1.00&=\underbrace{0.53}_{0}+\underbrace{0.25}_{1}+\underbrace{0.22}_{2}, &
0.53&=\underbrace{0.20}_{00}+\underbrace{0.18}_{01}+\underbrace{0.15}_{02},\\
0.25&=\underbrace{0.10}_{10}+\underbrace{0.08}_{11}+\underbrace{0.07}_{12}, &
0.07&=\underbrace{0.05}_{120}+\underbrace{0.02}_{121}.
\end{align*}
码长：$0.22\to1$；$0.20,0.18,0.15,0.10,0.08\to2$；$0.05,0.02\to3$，即 $(1,2,2,2,2,2,3,3)$。
\textbf{平均码长}（单位是三进制位 trit）：
\begin{align*}
\bar L_3&=0.22\times1+2\times(0.20+0.18+0.15+0.10+0.08)+3\times(0.05+0.02)\\
&=0.22+2\times0.71+3\times0.07=0.22+1.42+0.21=1.85 .
\end{align*}
\textbf{熵与界限}：$\Hx{X}=2.7535$ bit，$\lb3\approx1.5850$，
$\dfrac{\Hx{X}}{\lb3}=\dfrac{2.7535}{1.5850}\approx1.7373\le\bar L_3=1.85<2.7373$ $\checkmark$
效率 $\eta=1.7373/1.85\approx93.91\%$。\textbf{注意}：三元码每符号位数看着少（$1.85<2.72$），
但一个 trit 的信息量比一个 bit 多，\textbf{不能直接比大小}，要像上面除以 $\lb r$ 换成比特单位再比。
\end{calc}

\begin{qa}{为什么 Huffman 码不唯一？考试怎么答？}
三个原因，任一条都足以让答案不唯一：
(1) \textbf{左右可互换}——把 $0$ 和 $1$ 对调，码表全变但码长一点没变，$\bar L$ 不变；
(2) \textbf{概率相等时可随便排}——出现两个相等概率时谁左谁右无所谓，码字不同、码长相同；
(3) \textbf{中间序列的排序写法不唯一}——只要守住"取最小两个"的规则，插回位置可以有不同的写法。
标准答法：\textbf{最优码长分配 $(l_1,\dots,l_q)$ 唯一（去掉相等概率的对称性后），
但具体码字不唯一；所有 Huffman 码的平均码长都相同，都等于最优值。}
\end{qa}

%==============================================================
\subsection{其他几种常用变长编码（一句话认识）}

\begin{center}
\small
\begin{tabular}{@{}p{2.7cm}p{6.6cm}p{5.6cm}@{}}
\toprule
方法 & 做法（一句话） & 平均码长性质（一句话）\\
\midrule
Shannon 码 &
概率降序，取 $l_k=\left\lceil\lb\frac{1}{p_k}\right\rceil$，再用累积概率的二进制展开截 $l_k$ 位当码字。 &
即时码，但\textbf{不一定最优}；$\bar L<\Hx{X}+1$，$p$ 极端时浪费大
（如 $p=(0.9999,0.0001)$ 码长为 $1$ 和 $14$）。\\
\addlinespace
Fano 码 &
把符号\textbf{分成两组}使两组概率和尽量接近，一组记 $0$、一组记 $1$，再对每组递归地这样分。 &
不一定最优，只能保证 $\bar L\le\Hx{X}+2$ 量级，比 Huffman 略差。\\
\addlinespace
Shannon--Fano--Elias 码 &
用\textbf{累积概率中点} $\bar F(k)=F(k-1)+\frac{p_k}{2}$ 的二进制展开，
取 $\left\lceil\lb\frac{1}{p_k}\right\rceil+1$ 位当码字。 &
即时码，$\bar L<\Hx{X}+2$；价值在于\textbf{可递推实现}，是算术编码的前身。\\
\addlinespace
算术编码 &
把整条消息映射成 $[0,1)$ 里的一段\textbf{小区间}，用区间内某数的二进制展开当码字；
消息越长区间越窄，需要的位数自动跟着变。 &
\textbf{能逼近极限}：$N\to\infty$ 时冗余度趋于 $0$，极限下实现理想压缩，
这是它强于 Huffman 之处。\\
\addlinespace
游程编码（RLE） &
把二元序列中连续相同的符号换成"长度记数"，如 $00101110010001\cdots\to3,1,1,3,2,1,3,\cdots$
（"0 游程 3 个、1 游程 1 个……"）。 &
一一对应、可逆；本身是变长码，游程长度还需二次编码（Huffman 等），
对\textbf{长串重复}压缩效果极好（传真 MH 码就这么干）。\\
\bottomrule
\end{tabular}
\end{center}

\textbf{怎么记}：Shannon、Fano、SFE 都是"历史方案，都比 Huffman 差一点"；
SFE 的贡献是"递推 + 累积概率"的思想，被算术编码继承；游程编码是另一个赛道（先换表示，再压）。

%==============================================================
\subsection{平稳信源与马尔可夫信源（浅讲，够应付考试）}

\subsubsection{为什么还要熵率}
前面结论都建立在"符号之间独立"（无记忆）上。但真实的中文、英文显然不独立：
看到"信"，下一个字很可能是"息"、是"号"，不可能是"狗"——\textbf{前后文本身就带着信息}。
有记忆时每来一个符号带来的\textbf{新}信息比"独立时"少，因为一部分已被前文猜到了。
所以要一个新指标：\textbf{平均每个符号真正带来了多少新信息}。

\begin{defbox}{熵率与冗余度（4.4.1 节，必背）}
长度 $N$ 的符号块 $X_1X_2\cdots X_N$ 的联合熵为 $\Hx{X_1X_2\cdots X_N}$，
\textbf{平均每符号的熵}为 $H_N(X)=\frac1N\Hx{X_1X_2\cdots X_N}$，\textbf{熵率}定义为
\begin{align*}
H_\infty(X)=\lim_{N\to\infty}H_N(X)=\lim_{N\to\infty}\frac1N\Hx{X_1X_2\cdots X_N},
\end{align*}
而对平稳信源它还等于 $H_\infty(X)=\lim_{N\to\infty}\Hx{X_N\mid X_1\cdots X_{N-1}}$。
\textbf{冗余度}与\textbf{相对冗余度}（$K$ 为源字母表大小）：
\begin{align*}
\text{冗余度}=\lb K-H_\infty(X),\qquad
\text{相对冗余度}=1-\frac{H_\infty(X)}{\lb K}.
\end{align*}
\end{defbox}

\textbf{人话}：$H_\infty$ 就是"\textbf{把前后文都考虑进去以后，每个符号还剩多少真·新信息}"。
它一定不超过 $\Hx{X_1}$（独立时的熵）：\textbf{知道上下文只会减少不确定性}。
定理 4.4.1 把这个直觉说成三件事：$H_N(X)$ 随 $N$ 增大\textbf{单调不增}；
$H_N(X)$ 介于 $0$ 与 $\Hx{X_1}$ 之间；极限存在且有上面两个等价算法。
\textbf{冗余度}的人话：$\lb K$ 是"每符号最多能装多少比特"，$H_\infty$ 是"它实际装了多少"，
差出来的就是\textbf{浪费掉的部分}——也就是"这句话里有多少是语法和用词习惯替你写好的"。

\begin{calc}{英文的冗余度（教材例 4.4.1，背结论）}
英文 $26$ 个字母加空格共 $27$ 个字符，等概率时
$H_0=\lb27=\lb3^3=3\lb3=3\times1.5850=4.7549$ bit/字母。
只考虑字母频率得 $H_1=4.036$；考虑 $2$ 阶、$3$ 阶依赖得 $H_2=3.326$、$H_3=3.31$；
一般认为 $H_\infty\approx1.4$ bit/字母。于是
$\dfrac{H_\infty}{H_0}=\dfrac{1.4}{4.7549}\approx0.294$，$R=1-0.294\approx0.706=70.6\%$。
\textbf{读法}：写英文时约 $71\%$ 的内容是被语言结构"预先规定"的，只有约 $29\%$ 是作者自由选的。
\end{calc}

\subsubsection{马尔可夫信源：状态 + 转移矩阵 + 平稳分布}
\textbf{什么是马尔可夫信源}：信源任一时刻处于某个\textbf{状态} $s_j$，在 $s_j$ 下按
$\p{x_k\mid s_j}$ 发出符号 $x_k$，发完按 $p(s_j\mid s_i)$ \textbf{转移}到新状态。
关键特征是\textbf{下个状态和输出只取决于"现在处于哪个状态"，跟更早的历史无关}
（马尔可夫性/无后效性）。$m$ 阶马尔可夫信源（输出只与前 $m$ 个符号有关）的状态数是 $K^m$。

\textbf{状态转移矩阵} $P=(p_{ij})_{J\times J}$，$p_{ij}=p(s_j\mid s_i)$ 表示"从 $i$ 跳到 $j$"。
按概率的规矩，\textbf{每一行加起来必须等于 $1$}（从某状态出发总要跳到某处去）。
作业 4.12(1) 被点名"转移矩阵写错"，常见两种错：\textbf{行列搞反}
（$p_{ij}$ 是"$i\to j$"，\textbf{行=出发、列=到达}；搞反后"每行和为 $1$"变成"每列和为 $1$"，
后面平稳分布方程整体错位）；\textbf{把"该状态下输出符号的概率"当成转移概率}
（这是两套数）。\textbf{自查办法：填完矩阵立刻检查每行和是否为 $1$。}

\begin{defbox}{平稳分布怎么求（考试标准套路）}
稳态状态概率 $p_1,\dots,p_J$ 必须满足
\begin{align*}
\sum_{i=1}^{J}p_i\,p_{ij}=p_j\ (j=1,\dots,J),\qquad \sum_{j=1}^{J}p_j=1,
\end{align*}
矩阵写法即 $P^{\mathsf T}\boldsymbol p=\boldsymbol p$。
人话：稳态就是"转移一轮之后各状态的分布不再变"——\textbf{进这个状态的概率流正好等于出它的概率流}。
注意这 $J$ 个方程\textbf{线性相关}（只有 $J-1$ 个独立），
所以\textbf{必须配上 $\sum_jp_j=1$} 才解得出唯一答案，这一步千万别漏。
\end{defbox}

\begin{calc}{两状态小例子：手解平稳分布 + 算熵率}
两状态 $s_1,s_2$，转移矩阵（行=出发，列=到达）
$P=\begin{pmatrix}\frac34&\frac14\\[1pt]\frac12&\frac12\end{pmatrix}$。
先自查：第 1 行 $\frac34+\frac14=1$ $\checkmark$，第 2 行 $\frac12+\frac12=1$ $\checkmark$。

\textbf{(1) 解平稳分布}。设稳态概率 $p_1,p_2$，按定义列方程：
\begin{align*}
p_1&=\tfrac34p_1+\tfrac12p_2,\qquad
p_2=\tfrac14p_1+\tfrac12p_2,\qquad
p_1+p_2=1 .
\end{align*}
用第二式（最简单）：$p_2-\frac12p_2=\frac14p_1\Rightarrow\frac12p_2=\frac14p_1\Rightarrow p_2=\frac12p_1$。
代进 $p_1+p_2=1$：$p_1+\frac12p_1=1\Rightarrow\frac32p_1=1\Rightarrow p_1=\frac23,\ p_2=\frac13$。
\textbf{验算}（一定做）：$\frac34p_1+\frac12p_2=\frac34\cdot\frac23+\frac12\cdot\frac13=\frac12+\frac16=\frac46=\frac23$ $\checkmark$。

\textbf{(2) 各状态的条件熵}。设 $s_1$ 下发出 $0$ 的概率 $\frac34$、发出 $1$ 的概率 $\frac14$；
$s_2$ 下发出 $0,1$ 各 $\frac12$。则
\begin{align*}
\Hx{X\mid s_1}&=-\tfrac34\lb\tfrac34-\tfrac14\lb\tfrac14
=0.75\times0.4150+0.25\times2=0.3113+0.5=0.8113,\\
\Hx{X\mid s_2}&=-\tfrac12\lb\tfrac12-\tfrac12\lb\tfrac12=1 .
\end{align*}
（用到 $\lb\frac{4}{3}=\lb1.3333\approx0.4150$、$\lb4=2$。）

\textbf{(3) 熵率}（稳态下按状态概率加权）：
\begin{align*}
H_\infty(X)=\sum_{j=1}^{2}p(s_j)\Hx{X\mid s_j}
=\tfrac23\times0.8113+\tfrac13\times1=0.5409+0.3333\approx0.8742\ \text{bit/符号}.
\end{align*}

\textbf{(4) 冗余度}。源字母表 $\{0,1\}$，$K=2$、$\lb K=1$，
故冗余度 $=1-0.8742=0.1258$ bit，相对冗余度 $=1-0.8742/1=12.58\%$。

\textbf{(5) 与逐符号编码对比}。稳态下 $0,1$ 的边缘概率为
$\p{0}=\tfrac23\cdot\tfrac34+\tfrac13\cdot\tfrac12=\tfrac12+\tfrac16=\tfrac23$，$\p{1}=\tfrac13$。
对 $(\frac23,\frac13)$ 直接做二元 Huffman 只能得到两个长度 $1$ 的码字，$\bar L=1$ bit/符号，
效率 $0.8742/1=87.42\%$。这正好说明：\textbf{熵率 $0.8742$ 靠逐符号编码达不到，
必须把若干符号合起来编（扩展编码）才能逼近}——这就是下面平稳信源编码定理的意思。
\end{calc}

\begin{defbox}{平稳/马尔可夫信源的编码定理（结论记住即可）}
对离散平稳信源作 $r$ 元（教材记 $D$ 元）变长编码，把 $L$ 个源符号当一块编、
每符号平均码长记 $\bar l$，则 $\forall\varepsilon>0$，只要 $L$ 足够大，就存在唯一可译码使
\begin{align*}
\frac{H_\infty(X)}{\lb D}\ \le\ \bar l\ <\ \frac{H_\infty(X)}{\lb D}+\varepsilon .
\end{align*}
\textbf{翻译}：块越长，每符号平均码长就越能\textbf{任意逼近}熵率除以 $\lb D$。
无记忆时 $H_\infty(X)=\Hx{X_1}$，就退回第 5 节的 $\Hx{X}\le\bar L<\Hx{X}+1$。
\end{defbox}

%==============================================================
\begin{exam}{考试怎么考（第4章必考清单）}
\textbf{最高频的五类题}：
(1) \textbf{Huffman 全过程 + 平均码长 + 效率}（大题，$\sim10$ 分）。答题七步：
概率降序列表 $\to$ \textbf{写出每一轮合并过程}（$n-1$ 行）$\to$ \textbf{画码树}并在叶上标符号写码字
$\to$ 列码表（含 $l_i$、$p_il_i$）$\to$ $\bar L=\sum p_il_i$ $\to$ $\Hx{X}=\sum p_i\lb\frac1{p_i}$
$\to$ $\eta=\Hx{X}/\bar L$ 并补一句 $\Hx{X}\le\bar L<\Hx{X}+1$。
\textbf{不画树、只给码表要扣分}（作业 4.4 就栽在这里）。
(2) \textbf{Kraft 不等式判断}：给一组码长问能否构成 $r$ 元前缀码——写出通分后的和、与 $1$ 比较、
下结论；等于 $1$ 补"满树码"。
(3) \textbf{给码表判前缀码/唯一可译码}：$4$ 个小码表逐条判断并各算 $\bar L$；
按"非奇异 $\subset$ 唯一可译 $\subset$ 前缀"的层次说清结论。
(4) \textbf{求熵率与冗余度}：先写平稳分布方程（$P^{\mathsf T}\boldsymbol p=\boldsymbol p$ 加 $\sum p_j=1$）
解出 $p_j$，再算 $H_\infty=\sum_jp(s_j)\Hx{X\mid s_j}$，最后 $R=\lb K-H_\infty$；
\textbf{别忘了检查转移矩阵每行和为 $1$}。
(5) \textbf{概念小问}：$\bar L$ 的定义、三种码的包含关系、熵为什么是下限、
Huffman 码为何不唯一、三元码首步合并个数 $s=1+[(n-1)\bmod(r-1)]$。

\textbf{易丢分细节}：三元（$r$ 元）编码别忘\textbf{先算 $s$}，直接每轮合 $3$ 个会剩两个、树画不出来；
平均码长写清单位（bit/符号 或 trit/符号），效率写成百分数；
熵的 $\lb$ 是以 $2$ 为底（不是 $\ln$、不是 $\lg$）；
算 $\eta$ 时分子是熵、分母是平均码长，\textbf{别写反}（$\eta\le1$ 是天然检查）。
\end{exam}

\begin{qa}{为什么熵是平均码长的下限？}
分两层答。\textbf{(1) 数数}：$N$ 个符号的消息共 $q^N$ 种，要全部区分开至少需 $q^N$ 个不同码字，
即至少 $\lb q^N=N\lb q$ 位、每符号至少 $\lb q$ 位——正是"等概率时熵 $=\lb q$"。
\textbf{(2) 精确版}：设唯一可译码码长为 $l_i$，令 $q_i=\dfrac{r^{-l_i}}{\sum_jr^{-l_j}}$
（把 Kraft 的份额归一化成概率）。由 Kraft 不等式 $\sum_jr^{-l_j}\le1$ 得 $r^{-l_i}\le q_i$，
取对数得 $\lb\frac{1}{q_i}\le l_i\lb r$。于是
\begin{align*}
\sum_ip_i\lb\frac{1}{p_i}-\bar L\lb r
&\le\sum_ip_i\lb\frac{1}{p_i}-\sum_ip_i\lb\frac{1}{q_i}
=\sum_ip_i\lb\frac{q_i}{p_i}=-D(p\Vert q)\le0,
\end{align*}
最后一步用\textbf{相对熵（KL 散度）$D(p\Vert q)\ge0$}。故 $\bar L\ge\Hx{X}/\lb r$；
等号成立当且仅当 $q_i=p_i$，即 $l_i=\lb\frac{1}{p_i}$ 全是整数（概率都是 $2$ 的负整数次幂）。
\textbf{一句话}：熵是"消息本来需要的比特数"，码长是"你实际花掉的比特数"，
Kraft 不等式卡住了你能花的总额度，所以怎么编都不可能比熵更省。
\end{qa}

\begin{qa}{Huffman 为什么一定要"每次合并最小的两个"？合最大的会怎样？}
\textbf{合最小的两个}，等于让最少出现的两个符号共享同一段长路径，
它们占用的加权代价 $p_i\times(\text{多走的位数})$ 最小，对 $\bar L$ 伤害最小。
\textbf{反过来合最大的两个}，就等于给最常见的符号最长码字。拿第 6 节的分布真试一遍：
先合 $0.20$ 与 $0.19$ 得 $0.39$，再合 $0.39$ 与 $0.18$ 得 $0.57$……一直"合最大的"下去，
得到码长 $(6,6,5,4,3,2,1)$——\textbf{恰好倒过来了}，最常见的 $s_1$ 反而要 $6$ 位。
此时
$\bar L'=1.20+1.14+0.90+0.68+0.45+0.20+0.01=4.58$，
效率只剩 $2.6087/4.58\approx57\%$。所以"取最小两个"不是随便规定的，
它是能证明最优（定理 4.3.3）的关键一步。
（顺带一句：\textbf{"大概率配短码"只是必要条件，不等于 Huffman}。比如把码长配成 $(2,2,3,3,4,4,4)$，
虽然大的短小的长、Kraft 也过关（$\frac{2}{4}+\frac{2}{8}+\frac{3}{16}=\frac{15}{16}\le1$），
但 $\bar L=2.87$、$\eta\approx90.9\%$，明显不如 Huffman 的 $2.72$、$95.9\%$。）
\end{qa}
